\documentclass[11pt]{article}
\usepackage[T1]{fontenc}
\usepackage{lmodern}
\usepackage[margin=1in]{geometry}
\usepackage{amsmath,amssymb,amsthm,booktabs,longtable,array}
\usepackage[hidelinks]{hyperref}
\usepackage{xurl}
\setlength{\emergencystretch}{1em}
\newcommand{\IPRA}{\mathrm{iPRA}}
\newcommand{\PRA}{\mathrm{PRA}}
\newcommand{\PA}{\mathrm{PA}}
\newcommand{\IS}{\mathrm{I}\Sigma_1}
\newcommand{\Con}{\operatorname{Con}}
\newcommand{\Halt}{\operatorname{Halt}}
\newtheorem{proposition}{Proposition}
\title{Independent calibration and architecture review\\
\large PA194ORMP and PA194ORMPLC3}
\author{Research review, 15 November 2026}
\date{}
\begin{document}
\maketitle

\begin{abstract}
The two new 194-state machines preserve the resolved primitive program of
the revised AX11H construction.  Their consistency calibration can therefore
be inherited by an explicit, uniform source-to-tape simulation.  The
inspected logical interfaces of the revised calibration are coherent:
the current Foundation PA predicate is translated rather than identified
with a historical calculus, the forward proof suppliers are primitive
recursive, and the intuitionistic conclusion uses negative translation
without Markov's principle.  This review does not supply a completed Lean
formalization, rederive every inherited proof recipe, or establish a
machine with fewer than 194 states.  It also explains which ideas from the
ZFC, Diophantine, and concise-iPRA projects can be reused and which missing
steps prevent those projects from immediately yielding a smaller machine.
\end{abstract}

\section{Exact scope and imported results}

All repository references in this note are to ProveIt commit
\begin{center}\small
\path{7c0bb454cf73c29ca1c385f9fec856ddd236a7cb}.
\end{center}
Its Foundation gitlink is
\path{32e1a0956a8622fad067328ca1959729a7634428}, at
\path{lib/FormalizedFormalLogic-Foundation}.  That gitlink was checked
directly.  The mathematical parent is the \emph{revised} AX11H packet,
filed under
\begin{center}\small
\path{Computability/BusyBeaver/Research/PA/docs/pa194_ax11h_research_2026-11-14_2/}.
\end{center}
Its retained archive has SHA-256
\begin{center}\small
\path{93ee082ef035b2d32e71d6af5239b25b4bea1773f03bea4d2c39810a7fe15075}.
\end{center}
The suffix \path{_2} matters: an older packet with a similar name is a
different dependency.  The new construction retains the exact archive
as \path{calibration_AX11H.zip}, its ten proof fragments in
\path{calibration/}, and the exact calibrated table and source data in
\path{parent/}.

The imported written theorem is the final theorem of
\path{proof/consistency_equivalence.tex}: for the literal PA194AX11H
table, \(\IPRA\) proves its blank nonhalting equivalent to the pinned
PA consistency sentence.  This is an imported mathematical theorem,
not a hypothesis that PA is consistent.  Its packet explicitly states
that its finite execution checkers do not certify the entire arithmetic
derivation, and that it does not provide a new end-to-end Lean or Rocq
build.  Those qualifications also apply to the new machines' inherited
calibration.  Unfinished Lean work and the existence of a written
mathematical proof are different status claims.\cite{ax11h}

This review examined the complete final-equivalence and machine-simulation
fragments, the complete reverse calibration, the finite-demand and
stratum-ordering portions of the forward calibration, the current
Foundation adapter and its induction-code interface, and the actual
generic Lean one-way transfer statements.  It also examined the literal
relative checker and its receipts for both new machines.  It did
\emph{not} reconstruct every fixed numerical proof template or rebuild the
transitive Lean dependencies.  The accompanying source audit and symbolic
runtime replays provide the broader finite-data checks; they should not
be described as a newly completed proof-assistant kernel verification.

\section{The precise arithmetic endpoint}

The endpoint is
\path{DiophantineConsistency.IPRAPA196Consistency.nativePAConsistency}.
Its historical name contains ``196'', but its definition is a PA
consistency sentence, not the name of a particular machine.
It is the intuitionistic negation of the native reading of the mapped
Foundation PA-refutation predicate.  Three codings must be kept separate:
Foundation proof/formula codes, ordinary-PAeq proof/formula codes, and
the Cantor-coded numerical calculus used by the machine.  For example,
the adapter sends the historical false-formula code 13 to the numerical
contradiction 1; it does not identify those codes.\cite{adapter,transfer}

The file
\path{Computability/BusyBeaver/Research/PA/isigma1-machine/ConPAISigma1Directions.lean}
imports \path{ConPAUnconditional} and
\path{DiophantineConsistency.PRAPiOneImplication}.  It provides two
generic theorems.  With \(h(t)\) the chosen halting matrix, they require,
respectively, the following statements in every model of the specified
candidate PRA language whose arithmetic reduct satisfies \(\IS\):
\begin{align*}
 \text{mapped PA refutation}&\ \Longrightarrow\ \exists t\,h(t),\\
 \exists t\,h(t)&\ \Longrightarrow\ \text{mapped PA refutation}.
\end{align*}
The first endpoint additionally requires the matrix to be bounded; both
require it to be open.  Taking \(h(t)\equiv\chi_M(t)=1\), where
\(\chi_M\) is the primitive-recursive bounded simulation of the exact
literal table, meets those syntactic requirements because the matrix is
atomic.  The two endpoints then yield
\begin{align*}
 \IPRA&\vdash\neg\Halt(M)\to\Con(\PA),\\
 \IPRA&\vdash\Con(\PA)\to\neg\Halt(M).
\end{align*}
The actual Lean proof passes through classical model conservation and
negative translation.  It obtains double-negated positive existentials
where appropriate, then concludes the negative implications by
intuitionistic contraposition.  It does not eliminate double negation
from a positive existential and does not invoke Markov's principle.
The generic Lean source supplies these interfaces; it does not itself
instantiate them with the new literal machines.\cite{transfer}

\section{Why the written forward calibration is plausible in the weak base}

\subsection{The Foundation adapter}

The adapter translates between two arithmetic languages.  The source
has symbols \(0,1,+,\times,=,<\); the target has \(0,S,+,\times,=\).
It sends \(1\) to \(S0\) and \(s<t\) to an appropriately shifted
\(\exists z\,(s+Sz=t)\).  The written construction addresses bound-variable
shifts, simultaneous substitutions, universal closure, and the explicit
induction recognizer.  It does not appeal merely to the familiar
external equivalence of the two presentations.\cite{adapter}

A concrete version difference is handled: the historical finite PA axiom
\[
 x<y\to\exists z\,(z<y+1\mathbin{\wedge}x+z=y)
\]
and the current axiom
\[
 x<y\to\exists z\,(x+z=y)
\]
are not silently interchanged.  The adapter supplies the simple
existential projection in one direction and a fixed derivation using
the other finite arithmetic axioms in the other direction.  Likewise,
the current induction recognizer has bounded witnesses for the code of
the universal closure and its body; the code of an opened body is
bounded by the explicit opening function, not incorrectly by the
unopened formula code.  These details remove two possible failures of
an otherwise informal coding argument.\cite{adapter}

\subsection{Finite-demand compilation}

The forward numerical compiler cannot justify its construction by
quantifier-free induction on the unbounded assertion ``for every
environment, there exists a derivation.''  The revised proof instead
uses finite demand packets.  A query has a natural rank and a finite
parameter packet.  Its child queries are calculated before any child
proof is evaluated.  At most two child queries belong to the same
stratum, each with strictly smaller rank.  A complete binary tree of
depth bounded by the initial rank therefore contains every required
same-stratum query.  Inactive positions are filled with explicit
defaults, and the tree is evaluated in reverse order.  Its bounded
Boolean validity predicate can be maintained by quantifier-free
induction.\cite{forward,termleaf}

The essential restriction is that a call to a lower stratum may use a
larger argument, but no lower-stratum construction calls back into the
unfinished higher stratum.  That ordering is stated explicitly.
Term-value syntax, substitutivity, finite-member suppliers, and PA-leaf
suppliers are presented with finite witness tables or explicit
recurrences instead of an unbounded choice of a proof.  As a general
compilation method, this is a legitimate route to primitive-recursive
proof suppliers.  The present review checked these recursion interfaces
and their dependency order; the many fixed proof recipes they invoke
remain part of the imported written proof, rather than a newly replayed
kernel derivation.\cite{forward,termleaf}

\subsection{Cofinal enumeration and the carried store}

For a numerical derivation code \(e\) of contradiction, the command
compiler produces an \emph{all-store} halting word
\(V_e=A_0;W_e;\mathsf{Pop}\).  The final calibration wraps it in a
candidate whose raw code is
\[
 N_e=p\bigl(2,p(0,p(0,p(0,e_{\rm word}(V_e))))\bigr)\geq3.
\]
Here \(e_{\rm word}\) codes a command word, not a derivation.
Every earlier candidate drains within its seed plus one commands unless
halting has already occurred.  Thus
\(K(e)=(N_e+2)^2\) bounds the scheduler's halting time in abstract
commands.  Because \(V_e\) works from every store, the enumeration need
not reset the carried store or find the least successful candidate.
The exact source and physical clocks then turn this bound into a
primitive-recursive literal halting witness.\cite{ax11h,commands}

The reverse direction has a different obligation.  Its invariant says
that the formulas actually stored have internal PA proof witnesses.
It uses a finite vector of those witnesses, producing a
\(\Sigma_1\) invariant, and checks the guarded gate as well as the six
heads, the restricted \(\rho\) rule, and the diagonal rule.  Consequently,
the all-store forward property does \emph{not} authorize an arbitrary
initial source store for the reverse argument.  This distinction is
especially relevant to future attempts to reuse a 193-state runtime
with a different initializer.\cite{reverse}

\section{Uniform transfer to the new literal tables}

The new candidates are checked directly against the revised AX11H
parent.  Their SHA-256 bindings are:
\begin{description}
\item[PA194AX11H parent] \hfill\\[-1ex]
\path{c5dfe1026d077d3e062eca11edb789597bfbbced4f6040f8828ae5d467b1167a}
\item[PA194ORMP] \hfill\\[-1ex]
\path{a0473681d6d5fb68e5fccfc473f35bc67aa48083d847e349ba3c52865f293629}
\item[PA194ORMPLC3] \hfill\\[-1ex]
\path{6d8a8cf2f5578ccdc28dd7025bc41a7c279dfbc09c5d3893cb965d70292d10f4}
\end{description}
These hashes refer to the dense \path{machine.tm} bytes.  The archived
parent hash and both new hashes were recomputed during this review.
All three resolved instruction arrays have the same bytes, with SHA-256
\begin{center}\small
\path{43a4ba0d62f95fe766d5369940b0aa8b7a4da04824feb8836905e289640f59f7}.
\end{center}
This equality includes resolved control-flow targets, not only the
ordered sequence of primitive operations.  PA194ORMPLC3 retags 61
\path{continueTo} depths from 2 to 3 while retaining their resolved
targets.  A comparison of operation names alone would miss this
obligation.

The direct relative check covers all 3,671 primitive sites and all
4,778 guarded outcomes.  Exactly 4,777 outcomes are nonhalting.  It
checks padded destination addresses, absence of unintended halt rows,
nonnegative affine duration coefficients, and the unique halt outcome.
The four list-traversal schemas have separate one-block traces; list
induction extends them to arbitrary nonempty finite spare lists.
Registers and the gap parameter range over all natural numbers.
The full blank startup, including the zero exterior, reaches the common
source state at site 25 with the same logical data.  Startup durations
are 57,922, 58,014, and 58,010 for parent, ORMP, and ORMPLC3 respectively.
These are finite arithmetic bindings to the accompanying uniform
simulation proof, not a finite execution prefix offered as evidence of
nonhalting.

\begin{proposition}[The clock step needed for all-run transfer]
Suppose the source step is primitive recursive with an absorbing halt,
the literal machine has an absorbing semantic halt, and the frame
simulation is exact at cumulative times
\[
 T(k)=s+\tau(k),\qquad
 \tau(0)=0,\qquad
 \tau(k+1)=\tau(k)+D(z_k).
\]
Assume \(D(z_k)\geq3\) when the source state \(z_k\) is live, and
\(D(z_k)=0\) after source halt.  Then source halting at step \(k\)
implies literal halting by \(T(k)\).  Conversely, literal halting at
time \(t\) implies source halting by step \(t+1\).
\end{proposition}

\begin{proof}
The first claim is the exact endpoint simulation.  For the second,
suppose the source were still live at step \(t+1\).  Absorption makes
all earlier source states live, so
\(T(t+1)\geq s+3(t+1)>t\).  Exact simulation says the literal machine
is live at that later endpoint, contradicting literal halt at \(t\)
and its absorption.  This uses the duration lower bound only on a live
prefix; the clock need not keep growing after source halt.
\end{proof}

In the common blank-run invariant, gap is at most \(k\), the register
sum is at most \(2+k\), and the spare-list width is at most \(78+3k\).
The checked coefficient envelope gives
\[
 D(z_k)\leq698+18k,\qquad
 \tau(k)\leq9k^2+689k.
\]
The exact clocks are primitive recursive; the displayed polynomials
are convenient majorants.  Combining the exact source equality with
the endpoint simulations gives a primitive-recursive halting-witness
transport in both directions.  Consequently the same two model
implications, and hence the same negative \(\IPRA\) equivalence, apply
to both new machines.  No universal all-tape equivalence between their
raw transition graphs is required or claimed.  It is their proved
canonical-frame dynamics and full blank initialization that matter.

\section{What the three ancillary projects contribute}

\subsection{ZFC machines: useful physical methods, a different endpoint}

The pinned K215 project supplies a concrete 215-state blank-start binary
machine, including initialization, with 202 source subtrees, 4,751
primitive occurrences, and 6,176 runtime outcomes.  Its report records
compiled external Lean and Rocq equivalences with ZFC refutability, and
explicitly distinguishes those from an internal PA endpoint.  Its
95,714-step blank startup is part of the machine, not an auxiliary
unpriced preprocessing step.\cite{k215}

Several methods carry over: command-local field decoding, command
permutations, source reassociation, partial selectors, contextual
marker folds, and short history-dependent repairs.  In particular,
the K215 account explains a three-step replacement of the target of
row \(81/0\), freeing a reader row without adding a physical state.
The principle is to prove equality of complete local configurations
on the reachable histories, rather than require equality of the two
transition rows in isolation.  This is closely related to the useful
context repairs already present in the PA family.\cite{k215}

The K215 table itself cannot answer the present task: it is larger and
its endpoint is \(\Con(\mathrm{ZFC})\).  Replacing the target theory
requires a new calibrated source and literal pricing.  Likewise,
deleting an apparently unnecessary ZFC proof-theoretic operation does
not automatically delete a decoder state, because the operation's
reader may be shared by other source subtrees.

\subsection{Diophantine equations: arithmetic expression size is not state count}

The DiophantineConsistency project contains exact equations and several
native, standard, and equality-only \(\IPRA\) proof-existence endpoints.
Its principal equations encode ZFC consistency.  The project also
contains much of the PRA/\(\IS\)/iPRA transfer infrastructure used by
the PA calibration.  That infrastructure is immediately reusable as
proof support.\cite{diophantine}

A short equation is not yet a small blank-start machine.  A competing
design needs a PA-calibrated equation, an exhaustive tuple enumerator,
uniform integer arithmetic, a representation of its coefficients and
operations, and initialization.  The tuple decoder, polynomial
interpreter, sign handling, and literal coefficient data must all be
priced in the same two-symbol, L/R-only model.  A reduction in variable
count or arithmetic-operation count may increase those costs.

There is an informative PA-side example already in the compared
corpus.  The tagged-three recoding changes the leading formula radix
from four to three while retaining a uniform primitive-recursive
proof translation.  Its priced implementation reaches 207 states.
The reported optimizer score of 203 was not an accepted machine count,
because it overcounted four simultaneously compatible folds.  The
tagged-two prototype's optimistic 206 score fails a source contract.
Thus even a sound and apparently simpler numerical coding can be
larger at the literal-machine layer.\cite{tagged}

The Diophantine route remains open research, not refuted by these
observations.  A useful first experiment would price an entire fixed
polynomial evaluator and witness enumerator against the current
decoder, before investing in a new detailed consistency calibration.

\subsection{Concise iPRA: useful syntax, insufficient target strength}

The concise-iPRA project's stated objective is a theory bi-interpretable
with full first-order \(\IPRA\), with both numerical domains whole,
actual equality, and identity round trips.  Its smallest checker
metrics are normalized Python tokens, AST nodes, and bytes.  The native
scalar syntax uses variables, zero, successor, and an iterator;
formulas use equality and the two binders \(Q\) and \(D\).  Several
six- or seven-rule presentations are proved equivalent to the native
core.\cite{concise}

Those results provide promising normalization and iteration patterns.
They do not give a PA-refutation recognizer.  Enumerating contradictions
of an iPRA-equivalent calculus changes the target consistency sentence.
To use the syntax here, one would have to specify a PA-strength
extension, prove both syntactic refutation translations in the required
weak base, and implement its binding, substitution, normalization, and
proof validation.  Moving work from an inference rule into a canonical
syntax constructor may shorten a Python checker while enlarging the
binary interpreter.  Both the logical target and the literal price
therefore remain essential missing steps.

\section{Concrete directions below 194 and limits of the present search}

No new machine below 194 states was established in this review.  Several
seemingly easy logical ideas are already represented in the supplied
research: premise-order exchange in the six heads, alternative
diagonals, head and parameter permutations, H5X/J variants, different
formula radices, and substantial startup search.  Repeating them with
a different narrative would not supply a materially different design.
None of their negative finite search results proves a global lower
bound.

Two more focused directions remain technically meaningful.  First,
jointly synthesize source association and a local history repair.  The
relevant cost is the number of \emph{new} reader states after sharing
and compatible folds.  A reassociation that removes an old selector
but creates a new unshared subtree has saved nothing.  The most useful
certificate would list the complete predecessor histories of the
candidate deleted state, prove that each meets an existing reader at
an identical future configuration, and replay every affected symbolic
source contract.  The ORMP construction shows why source and runtime
cannot always be optimized independently; it does not itself show
that a further deletion is possible.

Second, a runtime with 193 states should be paired with a search over
\emph{provably admissible} initial scheduler configurations, not only
one prescribed initial tape.  Cofinal enumeration allows some variation
in the starting candidate counter, and the all-store forward compiler
provides useful flexibility.  However, reverse soundness still needs
explicit proof witnesses for every initial stored formula and a
correct scheduler phase.  An arbitrary canonical-looking register
frame at an arbitrary source address is insufficient.  A viable
startup family should therefore be specified first by a
primitive-recursive admissibility predicate, with proofs of source
soundness and cofinality for every member.  Only then can a literal
blank run that lands in that family finish the construction.

This second proposal is a research specification, not a claim that the
required startup exists.  An execution bound exhausted without a hit
leaves those executions unresolved; it cannot certify impossibility.
The two 194-state machines have a much firmer status: exact delivered
tables, uniform literal simulation evidence, and a clearly identified
written consistency calibration.  Their end-to-end Lean
formalization remains future work.

\begin{thebibliography}{12}
\bibitem{ax11h} Revised AX11H packet, pinned ProveIt path
\path{Computability/BusyBeaver/Research/PA/docs/pa194_ax11h_research_2026-11-14_2/};
especially \path{proof/consistency_equivalence.tex},
\path{proof/ipra_simulation_theorem.tex}, and
\path{proof/physical_ipra_audit.tex}.
\bibitem{adapter} Same packet,
\path{proof/foundation_paeq_completion.tex}, especially
``Binding the current Foundation predicate'' and
``Exact current Foundation endpoint.''
\bibitem{forward} Same packet,
\path{proof/forward_calibration_completion.tex}, finite-demand
construction and final stratum order.
\bibitem{termleaf} Same packet,
\path{proof/term_leaf_completion.tex}.
\bibitem{commands} Same packet,
\path{proof/source_command_completion.tex} and
\path{proof/ax11h_preservation.tex}.
\bibitem{reverse} Same packet,
\path{proof/consistency_implies_nonhalting.tex}.
\bibitem{transfer} Pinned ProveIt,
\path{Computability/BusyBeaver/Research/PA/isigma1-machine/ConPAISigma1Directions.lean},
theorems \path{ipra_nonhalt_implies_consistency} and
\path{ipra_consistency_implies_nonhalt}.
\bibitem{k215} Pinned ProveIt,
\path{Computability/BusyBeaver/Research/ZFC/candidate-k215/README.md}
and \path{proof.md}.  The proof-assistant build statements here are
reported project results; they were not freshly rebuilt for this review.
\bibitem{diophantine} Pinned ProveIt,
\path{Computability/DiophantineConsistency/README.md}.
\bibitem{concise} Pinned ProveIt,
\path{Logic/PrimitiveRecursiveArithmetic/ConciseIPRA/README.md},
``The question,'' ``The native scalar core,'' and the formalization table.
\bibitem{tagged} Pinned ProveIt,
\path{Computability/BusyBeaver/Research/PA/docs/pa194_lcfast_tagged_research_2026-11-14/README.md},
``Additional mathematical result and reduction attempts.''
\end{thebibliography}
\end{document}
